\documentclass{article}
\usepackage[utf8]{inputenc}
\usepackage[spanish]{babel}
\usepackage{geometry}
\geometry{margin=2.5cm}
\usepackage{amsmath,amssymb}
\usepackage{hyperref}

\begin{document}

\section{C4. Recomendación fundamentada del modelo de proceso}

Para cada uno de los cuatro escenarios analizados, la selección del modelo de proceso de requisitos se fundamenta en la evaluación de la estabilidad (K1), la disponibilidad de las partes interesadas (K2), la criticidad del sistema (K3), los riesgos principales (K4) y las exigencias normativas o contractuales (K5) \cite{Pohl2025}.

\subsection{Escenario 1: Extractora Río Lindo (Sistema Embebido)}
\begin{itemize}
    \item \textbf{Modelo recomendado:} Modelo Híbrido Controlado (Cascada e Iterativo/Incremental) \cite{GA_Prac_01}.
    \item \textbf{Desglose del modelo híbrido:} Se aplica un enfoque en \textbf{cascada} para los requisitos de seguridad críticos, límites de presión, temperatura y bloqueos definidos por el fabricante de los esterilizadores y exigidos por la aseguradora; y un enfoque \textbf{iterativo e incremental} para el desarrollo de la interfaz de operador, reportes y la integración con los sensores y válvulas aún pendientes de compra \cite{GA_Prac_01}.
    \item \textbf{Razones fundamentadas:}
        \begin{enumerate}
            \item \textit{K1 (Estabilidad):} Las reglas de seguridad y umbrales máximos son estables y fijas del fabricante, permitiendo congelar una línea base temprana \cite{Pohl2025}.
            \item \textit{K5 (Normativa):} La aseguradora exige formalmente un documento de requisitos de seguridad aprobado antes de la puesta en marcha \cite{GA_Prac_01}.
            \item \textit{K3 (Criticidad):} La alta criticidad por la seguridad física de los operadores exige una verificación formal estricta de las funciones de protección \cite{SWEBOK_v4_2024}.
        \end{enumerate}
    \item \textbf{Descarte de los otros modelos:} El modelo ágil puro se descarta porque la alta volatilidad en la interfaz no compensa el riesgo de desatender las reglas críticas de seguridad; el cascada puro resulta inflexible ante la incertidumbre de la selección de sensores en el tercer mes \cite{GA_Prac_01}.
    \item \textbf{Tensión principal:} La incertidumbre en la elección de los proveedores de sensores en el tercer mes frente a la exigencia de una especificación temprana de seguridad \cite{GA_Prac_01}.
\end{itemize}

\subsection{Escenario 2: Sabores de la Abuela Rosa (Aplicación Móvil)}
\begin{itemize}
    \item \textbf{Modelo recomendado:} Modelo Híbrido (Ágil e Iterativo Controlado) \cite{Inayat2015}.
    \item \textbf{Desglose del modelo híbrido:} Se aplica un enfoque \textbf{ágil} basado en un \emph{Product Backlog} vivo para la gestión del catálogo diario de menús, interfaz de usuario y promociones propuestas por la propietaria y su hijo; y un enfoque \textbf{iterativo controlado} para la pasarela de pagos por transferencia y los módulos de seguridad de datos personales conforme a la LOPDP \cite{GA_Prac_01}.
    \item \textbf{Razones fundamentadas:}
        \begin{enumerate}
            \item \textit{K1 (Volatilidad):} El menú diario cambiante y las ideas recurrentes de la dueña generan alta volatilidad que requiere repriorización continua \cite{GA_Prac_01, Pohl2025}.
            \item \textit{K2 (Disponibilidad):} La disponibilidad diaria de Doña Rosa para probar versiones maximiza la validación iterativa frecuente \cite{GA_Prac_01}.
            \item \textit{K4 (Plazo):} El plazo estricto de 14 semanas por las fiestas locales exige entregas incrementales de valor para mitigar el riesgo de descubrimiento tardío \cite{Sommerville2016}.
        \end{enumerate}
    \item \textbf{Descarte de los otros modelos:} El modelo en cascada se descarta por su rigidez ante los cambios diarios del negocio; el iterativo tradicional resulta excesivamente burocrático para un equipo reducido de dos desarrolladores a medio tiempo \cite{GA_Prac_01}.
    \item \textbf{Tensión principal:} La alta agilidad comercial requerida frente a la obligatoriedad estricta de cumplir con la protección legal de datos y conciliación de pagos \cite{GA_Prac_01}.
\end{itemize}

\subsection{Escenario 3: Sistema de Trámites Municipales (Contratación Pública)}
\begin{itemize}
    \item \textbf{Modelo recomendado:} Modelo Híbrido (Iterativo por Hitos Contractuales y Prototipado) \cite{GA_Prac_01}.
    \item \textbf{Desglose del modelo híbrido:} Se aplica un enfoque de \textbf{cascada por fases formales} para cumplir estrictamente con los tres entregables contractuales y de pago (mes 2, 6 y 7); y un enfoque \textbf{iterativo con prototipos} para investigar y resolver la integración técnica con la base de datos municipal antigua y no documentada \cite{GA_Prac_01}.
    \item \textbf{Razones fundamentadas:}
        \begin{enumerate}
            \item \textit{K5 (Contractual):} Los pliegos de contratación pública imponen fechas fijas, entregables aprobados y multas por cada día de retraso \cite{GA_Prac_01}.
            \item \textit{K4 (Riesgo técnico):} La integración con una base de datos heredada sin documentación requiere ciclos de investigación temprana \cite{GA_Prac_01, Pohl2025}.
            \item \textit{K1 (Cambios normativos):} Permite absorber de manera controlada la posible nueva ordenanza de tasas del cuarto mes sin romper la estructura de los entregables \cite{GA_Prac_01}.
        \end{enumerate}
    \item \textbf{Descarte de los otros modelos:} El modelo ágil puro se descarta por las restricciones legales, multas y la inflexibilidad de alcance de la contratación pública; el cascada rígido se descarta porque la ordenanza futura invalidaría una única especificación inicial \cite{GA_Prac_01}.
    \item \textbf{Tensión principal:} La rigidez normativa de los pliegos y multas de contratación pública frente a la incertidumbre técnica de la base de datos heredada \cite{GA_Prac_01}.
\end{itemize}

\subsection{Escenario 4: Cooperativa Frutos del Estero (Plataforma Web Rural)}
\begin{itemize}
    \item \textbf{Modelo recomendado:} Modelo Híbrido (Iterativo-Incremental y Ágil de Campo) \cite{GA_Prac_01}.
    \item \textbf{Desglose del modelo híbrido:} Se aplica un enfoque \textbf{iterativo e incremental} para el núcleo de cálculos de liquidación y créditos (reglas estables del estatuto), combinado con ciclos \textbf{ágiles y de prototipado offline} para el registro de entregas en los centros de acopio rurales con conectividad intermitente \cite{GA_Prac_01}.
    \item \textbf{Razones fundamentadas:}
        \begin{enumerate}
            \item \textit{K2 (Disponibilidad mixta):} Permite sincronizar las revisiones formales con la directiva mensual y la validación de los socios en asambleas trimestrales y días de acopio \cite{GA_Prac_01}.
            \item \textit{K3 (Criticidad técnica):} Facilita comprobar de manera anticipada el funcionamiento sin conexión y la posterior sincronización de datos \cite{GA_Prac_01}.
            \item \textit{K5 (Informes):} La generación de incrementos funcionales semestrales respalda directamente los informes exigidos por el donante \cite{GA_Prac_01}.
        \end{enumerate}
    \item \textbf{Descarte de los otros modelos:} El modelo ágil se descarta porque la directiva se reúne solo una vez al mes y los socios están dispersos, impidiendo la priorización continua; el cascada no permitiría validar oportunamente el funcionamiento offline \cite{GA_Prac_01}.
    \item \textbf{Tensión principal:} Las exigencias de reportes formales semestrales del donante frente a la baja frecuencia de retroalimentación directa de los productores rurales \cite{GA_Prac_01}.
\end{itemize}

\section{C5. Análisis de riesgos del modelo elegido y mitigación}

\begin{enumerate}
    \item \textbf{Riesgo 1 en Escenario 1 (Híbrido controlado: Cascada + Iterativo/Incremental):} Obsolescencia de los requisitos ante cambios en los sensores seleccionados en el tercer mes \cite{GA_Prac_01}.
        \begin{itemize}
            \item \textit{Impacto:} Afecta el diseño del hardware y genera retrabajo en la integración del sistema embebido.
            \item \textit{Mitigación:} Establecer una capa de abstracción de hardware (HAI) que aísle el software de control de la marca específica de sensores.
        \end{itemize}
    \item \textbf{Riesgo 2 en Escenario 1 (Híbrido controlado: Cascada + Iterativo/Incremental):} Dificultad de validación continua por la disponibilidad baja a media del jefe de planta (dos mañanas por semana) y turnos rotativos de operadores \cite{GA_Prac_01}.
        \begin{itemize}
            \item \textit{Impacto:} Retrasos en la aprobación formal de los flujos de control y detección tardía de malentendidos operativos.
            \item \textit{Mitigación:} Programar revisiones estructuradas fijas al cierre de cada incremento y validar mediante prototipos de interfaz antes de la codificación final.
        \end{itemize}

    \item \textbf{Riesgo 1 en Escenario 2 (Híbrido: Ágil + Iterativo controlado):} Omisión de Requisitos No Funcionales (RNF) de seguridad y cumplimiento de la LOPDP en la aplicación móvil \cite{GA_Prac_01}.
        \begin{itemize}
            \item \textit{Impacto:} Incumplimiento de la normativa legal ecuatoriana de protección de datos personales de los clientes.
            \item \textit{Mitigación:} Incorporar historias de usuario de seguridad y privacidad como criterios de aceptación obligatorios desde el primer sprint \cite{Pohl2025}.
        \end{itemize}
    \item \textbf{Riesgo 2 en Escenario 2 (Híbrido: Ágil + Iterativo controlado):} Expansión descontrolada del alcance (\emph{scope creep}) debido a las ideas recurrentes de la propietaria y su hijo en cada reunión \cite{GA_Prac_01}.
        \begin{itemize}
            \item \textit{Impacto:} Agotamiento prematuro del presupuesto cerrado y retraso en la entrega para las fiestas locales.
            \item \textit{Mitigación:} Gestionar rigurosamente el \emph{Product Backlog} aplicando priorización estricta (técnica MoSCoW) para congelar el alcance de cada incremento.
        \end{itemize}

    \item \textbf{Riesgo 1 en Escenario 3 (Híbrido: Cascada por hitos + Iterativo con prototipado):} Desviación del alcance acordado en el entregable del mes 2 debido a modificaciones en el flujo de aprobación del GAD \cite{GA_Prac_01}.
        \begin{itemize}
            \item \textit{Impacto:} Riesgo de retrasos y aplicación de multas contractuales estipuladas en los pliegos.
            \item \textit{Mitigación:} Gestionar los cambios mediante un comité formal de control de cambios (CCB) documentado en actas de iteración \cite{Pohl2025}.
        \end{itemize}
    \item \textbf{Riesgo 2 en Escenario 3 (Híbrido: Cascada por hitos + Iterativo con prototipado):} Incompatibilidad o errores de lectura en la base de datos municipal antigua y no documentada \cite{GA_Prac_01}.
        \begin{itemize}
            \item \textit{Impacto:} Fallas críticas en la consulta de datos catastrales durante la emisión de trámites.
            \item \textit{Mitigación:} Dedicar la primera iteración del contrato a realizar ingeniería inversa y pruebas de concepto (PoC) sobre la estructura de la base de datos heredada.
        \end{itemize}

    \item \textbf{Riesgo 1 en Escenario 4 (Híbrido: Iterativo-Incremental + Ágil de campo):} Fallas en la sincronización de registros de cacao debido a la inestabilidad de internet en el centro de acopio de Mocache \cite{GA_Prac_01}.
        \begin{itemize}
            \item \textit{Impacto:} Pérdida temporal de datos de entrega y descontento en los pagos a los productores.
            \item \textit{Mitigación:} Priorizar pruebas tempranas de almacenamiento local offline (IndexedDB/SQLite) y colas de sincronización robustas en los primeros incrementos.
        \end{itemize}
    \item \textbf{Riesgo 2 en Escenario 4 (Híbrido: Iterativo-Incremental + Ágil de campo):} Desalineación entre los requerimientos de la directiva (reunión mensual) y las necesidades reales de los socios dispersos (asambleas trimestrales) \cite{GA_Prac_01}.
        \begin{itemize}
            \item \textit{Impacto:} Rechazo de la plataforma por parte de los usuarios finales por no ajustarse a su operación en recintos rurales.
            \item \textit{Mitigación:} Nombrar comités de enlace durante los días de acopio para validar incrementos funcionales periódicos que respalden los informes semestrales del donante.
        \end{itemize}
\end{enumerate}

\begin{thebibliography}{9}
\bibitem{Pohl2025} Klaus Pohl. \emph{Requirements Engineering: Fundamentals, Principles, and Techniques}. Springer, 2nd edition, 2025.
\bibitem{SWEBOK_v4_2024} Hironori Washizaki, editor. \emph{Guide to the Software Engineering Body of Knowledge (SWEBOK Guide)}, Version 4.0. IEEE Computer Society, 2024.
\bibitem{Inayat2015} Irum Inayat et al. A systematic literature review on agile requirements engineering practices and challenges. \emph{Computers in Human Behavior}, 51:915-929, 2015.
\bibitem{Sommerville2016} Ian Sommerville. \emph{Software Engineering}. Pearson, 10th edition, 2016.
\bibitem{GA_Prac_01} Universidad Técnica Estatal de Quevedo (UTEQ). \emph{Material complementario GA-PRAC-01: Proceso de requisitos}, 2026.
\end{thebibliography}

\end{document}